\section*{Введение}
\addcontentsline{toc}{section}{Введение}

\textbf{Актуальность темы.}  
В современных условиях цифровизации сферы общественного питания возрастает потребность в создании удобных и функциональных веб-приложений, позволяющих автоматизировать процессы обслуживания клиентов в кафе и ресторанах. Традиционные способы оформления заказов, бронирования столиков и взаимодействия с персоналом зачастую являются разрозненными и не обеспечивают достаточного уровня удобства для пользователей и эффективности для заведений.

Особую актуальность данная проблема приобретает в условиях развивающихся рынков, где цифровая инфраструктура ресторанного бизнеса находится на начальном этапе формирования. В частности, в Туркменистане, включая столицу Ашхабад, рынок цифровых сервисов для кафе и ресторанов развит слабо, а существующие решения не обеспечивают комплексной автоматизации процессов заказа, бронирования и предварительного планирования посещений.

Разработка единой веб-платформы, объединяющей функции онлайн-заказа, бронирования столиков и предзаказа блюд, является актуальной задачей, способной повысить уровень сервиса в заведениях общественного питания и улучшить пользовательский опыт.

\textbf{Степень разработанности проблемы.}  
На международном рынке существуют крупные цифровые платформы, такие как сервисы доставки еды и системы бронирования столиков. Однако большинство из них ориентированы на развитые рынки и не учитывают специфику локальных экономических условий и инфраструктуры. В частности, в странах с ограниченным развитием безналичных платежей и цифровых банковских сервисов отсутствуют комплексные решения, адаптированные под смешанные способы оплаты, включая наличный расчет и оплату через мобильные устройства.

Таким образом, проблема создания универсального веб-приложения для автоматизации процессов в кафе и ресторанах остается недостаточно решенной, особенно в контексте локального рынка города Ашхабад.

\textbf{Объект исследования} — процессы обслуживания клиентов в заведениях общественного питания (кафе и ресторанах).

\textbf{Предмет исследования} — веб-технологии и программные средства автоматизации процессов заказа, бронирования столиков и предзаказа блюд.

\textbf{Целью данной работы} является разработка веб-приложения для управления заказами, бронированием столиков и предзаказом блюд в сети кафе, ориентированного на использование в условиях развивающегося рынка общественного питания города Ашхабад с возможностью дальнейшего масштабирования в формате стартапа.

\textbf{Практическая значимость.}  
Разрабатываемая система может быть использована в кафе и небольших сетях общественного питания для повышения уровня автоматизации обслуживания клиентов, сокращения времени ожидания заказов и улучшения качества сервиса. Особое значение проект имеет для рынка Туркменистана, где цифровизация ресторанного бизнеса находится на начальном этапе развития.

Дополнительно система учитывает особенности локальной инфраструктуры, включая ограниченное распространение банковских карт и цифровых платежных систем. В связи с этим предусматривается поддержка различных способов оплаты, включая наличный расчет и оплату через мобильные устройства, что повышает доступность приложения для конечных пользователей.

В перспективе проект рассматривается как основа для создания стартапа, ориентированного на внедрение цифровых решений в сети кафе города Ашхабад.